\begin{lemma}\label{lem: euclidean}
    For the Euclidean Algorithm, let $n$ be a natural number and $q$ be a positive natural number. 
    Then, there exist \textbf{unique} natural numbers $m$ and $r$ such that $0 \leq r < q$ and $n = mq + r$.
\end{lemma}
\begin{proof}
    For the existance, this has already been done beforehand in Chapter 2C, Problem 5.
    To prove uniqueness, suppose there exist two representations of $n$ as expressed by $n = m_1 q + r_1$ and $n = m_2 q + r_2,$
    where as $r_1, r_2 \in [0, q)$ (i.e., $0 \leq r_1, r_2 < q)$).

    Notice that if we subtract these two representations, we get $n - n= m_1 q + r_1 - (m_2 q + r_2) = (m_1 - m_2) q + (r_1 - r_2) = 0.$
    This implies that $(m_1 - m_2) q = r_2 - r_1.$ 
    Because $r_1, r_2 < q$, we have that $-q < r_2 - r_1 < q$.
    Substituting, we get that $-q < (m_1 - m_2) q < q$
    We can divide the inequalities to get $-1 < m_1 - m_2 < 1.$
    The only integer in between -1 and 1, exclusive, is 0.
    Hence, $m_1 - m_2 = 0$ implies $m_1 = m_2$.
    Moreover, this would also imply $(m_1 - m_2)q = 0q = 0 = r_2 - r_1.$ Then, $r_1 = r_2.$
    We have now shown that $m_1 = m_2$ and $r_1 = r_2.$ Therefore, $m$ and $q$ are unique. 


\end{proof}

\begin{problem}{Interspesing of integers by rationals}
    Let $x$ be a rational number. 
    Then, there exists an integer $n$ such that $n \leq x < n + 1.$
    In fact, this integer is unique (i.e., for each $x$, there is only one $n$ for which $n \leq x < n + 1)$.
    In particular, there exists a natural number $N$ such that $N > x$ (i.e., there is no such thing as a rational number which is larger than all the natural numbers.)
\end{problem}

\begin{proof}
    Let $x = x_1 \quot x_2$ be a rational number where $x_1$ is an integer and $x_2$ is a non-zero integer.
    We consider three cases by the law of trichotomy on order for the integers.
    \begin{enumerate}
        \item Say $x \geq 0.$ Then, $x_1$ is a natural number while $x_2$ is a positive natural number. 
        Then, by the Euclidean Algorithm, there exist unique natural numbers $m$ and $r$ such that $0 \leq r < x_2$ and $x_1 = m x_2 + r.$
        Notice that $mx_2 \leq m x_2 + r < (m + 1) x_2.$ Since $x_2$ is a non-zero integer, we can divide the inequality to get $m \leq x < m + 1$.
        \item Say $x < 0,$ so there exists a positive rational $x'$ such that $x = -x'.$ 
        Hence, we can use the previous case to determine that there exists an integer $m$ such that $m \leq x' < m + 1.$
        Multiply all sides by -1 to get $-m \geq -x' > -m - 1$ (i.e., $-m-1 < x \leq -m$). We now consider two sub-cases.
        \begin{enumerate}
            \item Say $x = -m.$ Then, take $n:= -m$, so that $n \leq x < n + 1.$
            \item Say instead that $x \neq -m$. Then, take $n := -m - 1$ (i.e., n + 1 = -m). Then, $-m-1 < x \leq -m$ can rewritten as $n < x \leq n +1$.
        \end{enumerate}
    \end{enumerate}
    In either case, there exists an integer $n$ such that $n \leq x < n +1.$
    To prove uniqueness, suppose there are two distinct $n \neq m$ integers such that $n \leq x < n+1$ and $m \leq x < m + 1.$
    Then, assume WLOG that $n < m$. If $m > n$ since we can just repeat what follows by swapping $n = m.$
    Notice $n < m \leq x < n +1$, which can be simplified into $n < m < n +1.$
    This implies that $m = n + d$ and $n + 1 = m + d'$ where $d$ and $d'$ are positive.
    Notice then that $n + 1 = n + d + d'$, which can be simplified into $1 = d + d'.$
    Since $d$ is a positive number, there exists a natural number $k$ such that $d = k +1.$ 
    Notice then $1 = k + 1 + d$ can be simplified into $0 = k + d.$
    Therefore, $d = 0.$ However, this contradicts that $d$ is positive!
    %with different representations of $x_1 \quot x_2 = \tilde{x}_1 \quot \tilde{x}_2$ where $x_1, \tilde{x}_2$ are integers and $x_2, \tilde{x}_2$ are non-zero integers.
    
    Finally, to show that there exists a natural number $N$ such that $N > x$, use the just proven problem that there exists an integer $n$ such that $n \leq x < n + 1.$
    Choose $N := n + 1.$ 
    Therefore, $N = n + 1 > x.$

\end{proof}

This part is just a reflection on how to do problem solving for this one since I actually got the idea of using the Euclidean algorithm. 
However, I was still severely lacking with the problem solving aspect because:
\begin{enumerate}
    \item I was genuinely just planning to use only the quotient alone. Like in the $n = mq + r$, I was only going to use $m$.
    This is because my examples were all so small that $m$ could have been used. What I learned is to use bigger and more varied examples so that I can check my assumptions.
    \item Another weakness is that I could not figure out how to deal with the negative case. 
    I was having trouble on how to deal with the fact that $-\frac{3}{2}$ is going to have its lower portion as $n=-2$ and not $n=-1.$
    But, the lower portion (the $n$) of $-\frac{4}{2}$ is $n=2$. Like, what the solution that was presented was smart in just divying up the problem into cases.
    Hence, don't be afraid to do cases.
    \item To prove the uniqueness, that was also hard because tbh, proving uniqueness is so hard for me. 
    There are usually two ways. Suppose you have two different solutions; show they are actually the same.
    Or, you can also do: suppose they are distinct; make this lead into a contradiction.
\end{enumerate}
\begin{definition}{Infinite Descent}
    A definition: a sequence $a_0, a_1, \dots$ of numbers is said to be in \textit{infinite descent} if we have $a_n > a_{n+1}$ for all natural numbers $n$ (i.e., $a_0 > a_1 > a_2 > \cdots$).
\end{definition}
\begin{problem}{Principle of Infinite Descent}
    Prove the principle of infinite descent.
\end{problem}
\begin{proposition}
    It is not possible to have a sequence of natural numbers which is in infinite descent.
    (Hint: assume for sake of contradiction that you can find a sequence of natural numbers which is in infinite descent.
    Since all the $a_n$ are natural numbers. you know that $a_n \geq 0$ for all $n$.
    Now use induction to show in fact that $a_n \geq k$ for all $k \in \mathbb{N}$ and all $n \in \mathbb{N}$, and obtain a contradiction.)
\end{proposition}
\begin{proof}
    Suppose for the sake of contradiction that there is a sequence $\big(a_n\big)$ that is in infinite descent.
    We will show that for all natural numbers $k$ and $n$, we have $a_n \geq k.$ We do this by inducting on $k.$
    For the base case of $k = 0$ notice that for all $n \in \mathbb{N}, \, a_n \geq 0$ because $a_n \in \mathbb{N}$ from Chapter 2, Section A, Problem 4, Proposition 1.

    Now, we will prove the inductive step. 
    Suppose that for any natural number $n$, $a_n \geq k$. Hence, $a_{n+1} \geq k$. 
    By the sequence being in infinite descent, we have that $a_n > a_{n+1}.$
    By the property of transitivity, $a_n > a_{n+1} \geq k$. Then, $a_n > k$ (i.e., $k < a_n.$)
    Therefore, $k + 1 \leq a_n$ (i.e., $a_n \geq k +1$).

    Now, we will show the contradiction. Let $k := a_0 + 1.$ 
    This turns $a_0 > k$ to $a_0 > a_0 + 1$, which finally becomes $0 > 1$, clearly a contradiction1
\end{proof}

Reflection for this one is to use the principle of transitivity and the power of the `for all' logic operator. 
I did not really appreciate the fact that we had the tool of the `for all $n \in \mathbb{N}$' with me, so I didn't notice that $a_{n + 1} \geq k.$
ALSO, another example where I did not appreciate the for all is when they let $k := a_0$ since the inductive hypothesis could work for any $k$.
Finally, the proof for this one is so weird. 
Like, genuinely, they assume something false (normal for a contradiction). Go on a long tangent of proving through induction, and finally doing a simple contradiction.
\begin{proposition}
    Does the principle of inifinite descent work if the sequence $a_1, a_2, a_3, \dots$ is allowed to take integer values instead of natural number values?
    What about if it is allowed to take positive rational values instead of natural numbers? Explain. 
\end{proposition}
\begin{proof}
    Yes, consider for the integers $a_n = -n$ and for the rationals $a_n = \frac{1}{n}$.
\end{proof}

\begin{problem}{Formalizing that there does not exist a $x \in \mathbb{Q}$ such that $x^2 = 2.$}
    We solve each proposition. Let $p, q$ be positive integers.
\end{problem}
\begin{proposition}
    Every natural number is either even or odd, but not both.
\end{proposition}
\begin{proof}
    Every natural number is either even or odd by the Euclidean algorithm with divisor 2 since it states that for every natural number $n$, $n = 2q + r$ where $q$ is an integer and $0 \leq r < 2.$
    Equivalently, $r = 0, 1.$

    To show a number cannot both be odd and even, suppose for the sake of contradiction that $n$ is even and odd.
    Then, $n = 2k$ and $n = 2l + 1$ for some natural numbers $k$ and $l$.
    We consider by trichotomy three cases:
    \begin{enumerate}
        \item Say $k < l$.
        Then, $n = 2k < 2l < 2l + 1$. Therefore, $n \neq 2l + 1$, a contradiction!
        \item Say $k = l$. Then, $n = 2k = 2l \neq 2l + 1$, a contradiction!
        \item Say $k > l$ (i.e., $l < k$). Then, $l + 1 \leq k$, so $2l + 2 \leq 2k$.
        By Proposition 2.2.12e ($a < b \iff a + 1 \leq b$), $2l + 1 < 2l + 2 \leq 2k = n.$ Therefore, $n \neq 2l + 1$, a contradiction!
    \end{enumerate}
    In all cases, we had a contradiction due to assuming $n$ was both even and odd.
    Therefore, a number $n$ can either be only even or odd.
\end{proof}
\begin{proposition}
    If $p$ is odd, then $p^2$ is also odd.
\end{proposition}
\begin{proof}
    Suppose $p$ is odd, so $p = 2k + 1$ for some integer $k$. 
    Then, $p^2 = (2k + 1)^2 = 4k^2 + 2k + 1 = 2(k^2 + k) + 1.$
    Therefore, $p^2$ is also odd.
\end{proof}

\begin{proposition}
    Since $p^2 = 2q^2$, we have $q < p$.
\end{proposition}
\begin{proof}
    Suppose for the sake of contradiction that $p^2 = 2q^2$ and $q \geq p.$
    Then, $qp \geq p^2 = 2q^2$. Dividing both sides by $q$ gives $p \geq 2q.$ 
    Combining the equalities gives us $2q \leq p \leq q.$ 
    Simplifying gives us $2q \leq q$ which implies $2 \leq 1$, a contradiction!
\end{proof}

